% Source-derived chapter 09: Function-theoretic corollaries
% Original source: intersection-complexes-unramified-l-factors
\chapter{Function-theoretic corollaries}
\label{intersection-complexes-unramified-l-factors-chapter-09}
\source{intersection-complexes-unramified-l-factors}

\begin{remark}[Source section]
\label{intersection-complexes-unramified-l-factors-chapter-09-source}
\source{intersection-complexes-unramified-l-factors}
\source{intersection-complexes-unramified-l-factors}
This chapter reproduces the corresponding section of the retrieved
LaTeX source, including its definitions, statements, examples, and
proofs. It has not yet been translated into Lean.
\notready
\end{remark}

% The source section title was: Function-theoretic corollaries
\subsection{Pushforward of the basic function} \label{sect:finitefields} 
\source{intersection-complexes-unramified-l-factors}

When $k$ is the algebraic closure of a finite field $\mathbb F$, and $X$ is defined over $\mathbb F$, satisfying the assumptions of \S \ref{assumptions-finitefield}, the action of the geometric Frobenius $\Fr$ morphism on $ \bar\pi_!(\IC_{\barY^{\ch\lambda}})$ is described, up to a yet unknown permutation action on the set $\mB^+_X = \bigcup_{\ch\lambda} \mB_{X,\ch\lambda}$ of central components of critical dimension, by Proposition \ref{prop:format}. We use this to prove Theorems \ref{theorem-affine-closure} and \ref{theorem-general} from the introduction.




\begin{proof}[Proof of Theorems \ref{theorem-affine-closure} and \ref{theorem-general}]
Recall that $\mf o$, in the context of these theorems, denotes the ring $\mathbb F\tbrac t$, where $\mathbb F$ is the finite field of definition of $X$, and $F$ is its fraction field. 

We need to recall the definition of the IC function $\Phi_0$  from \cite{BNS}: It is a function on $(X(\mf o)\cap X^\bullet(F))/G(\mf o)$ which, in our case, is parametrized by the set $(\mf c_X^-)^{\Fr}$ of elements of $\mf c_X^-$ that are fixed under the Galois group. To define it, choose an arc $\gamma$ in the coset of such an element $\ch\theta$, and consider a finite-dimensional formal model $\wh{Y}_y$ of the formal neighborhood of $\gamma$ in the arc space $\msf L^+X$ (Definition \ref{def:formalmodel}).  
In our case, we can take $Y= \sY^{\ch\theta}$ and $y=$ the point $t^{\ch\theta}$ on the central fiber $\msf Y^{\ch\theta,\ch\theta}$, by Theorem \ref{thm:DGK}. Then, the value of $\Phi_0$ on $\ch\theta$ is equal to the trace of geometric Frobenius on the stalk of the intersection complex $\IC_Y$ at $y$, \emph{where the intersection complex is normalized to be constant} (without Tate or cohomological twists) \emph{on the smooth locus of $Y$}. 

In our setting, this means that for every component $\mathscr Y$ as in Proposition \ref{prop:format}, the Tate and cohomological twist on $\IC_{\mathscr Y}$ should be modified from $\ol\bbQ_\ell(\frac{\dim\mathscr Y}{2})[\dim \mathscr Y]$ to $\ol\bbQ_\ell$, and \eqref{e:defVX} should be replaced by the space 
\begin{equation}\label{e:defVX2}
\source{intersection-complexes-unramified-l-factors}
\bigoplus_{\mathscr Y} \bigoplus_{\mB_{\mathscr Y}} \ol\bbQ_\ell(-\smallfrac{\dim\mathscr Y}{2})[-\dim \mathscr Y] = \bigoplus_{\msf b\in \mB_{X,\ch\lambda}} \ol\bbQ_\ell(-\dim \msf b - \smallfrac{1}{2})[-2\dim \msf b - 1].
\end{equation}


To calculate the value at $\ch\lambda(t)$ of the integral of the basic function that was denoted by $\pi_! \Phi_0$ in the introduction, for $\ch\lambda$ fixed by Frobenius, we need to calculate the (alternating) trace of Frobenius on the fiber of $ \bar\pi_!(\IC_{\barY^{\ch\lambda}})$ over an $\mathbb F$-point of the diagonal $C\hookrightarrow \sA^{\ch\lambda}$, and then divide by the factor $\tr_{\check T}(\Fr, \Sym^\bullet (\check{\mathfrak n}(1)))$ in order to account for the difference between $\IC_{\sY}$ and $\IC_{\barY}$ (Corollary \ref{cor:pibarY}).

Taking into account the twists in the intersection complexes of the $C^{\mf R}$'s in \eqref{e:format1}, we deduce that, with this normalization of the IC sheaves, Frobenius acts on that fiber as on
\[
\bigoplus_{\deg(\mf R)=\ch\lambda} 
\Bigl( \bigotimes_{\ch\mu} 
\Sym^{N_{\ch\mu}}(\bigoplus_{\msf b\in \mB_{X,\ch\mu}} \ol\bbQ_\ell(-\dim \msf b)[2\dim \msf b] )\Bigr).
\]

Thus, in the notation of Theorems \ref{theorem-affine-closure} and \ref{theorem-general}, we have
\[\pi_! \Phi_0 = \tr_{\check T}(\Fr, \Sym^\bullet (\check{\mathfrak n}(1)))^{-1} \cdot \tr_{\check T}(\Fr, \Sym^\bullet (\bigoplus_{\msf b\in \mB_X^+} \ol\bbQ_\ell(-\dim \msf b))),\]
where we remind that $\mB_X^+ = \bigcup_{\ch\mu\in \mf c_X} \mB_{X,\ch\mu}$. 

The dimensions $\dim \msf b$ are given by Proposition \ref{prop:Vbasis}:
they are either equal to $\frac 1 2(\len(D)-1)$, if $\msf b$ meets $\sY_{X^\bullet}^{D}$, or $\brac{\rho_G,\ch\lambda-\ch\theta}$, if $\msf b$ meets $\sY^{\ch\lambda,\ch\theta}$ for $\ch\theta \in \Cal D^G_{\mathrm{sat}}(X)$. 

Both theorems assume the existence of a $G$-eigen-volume form on $X^\bullet$, which we fix, with eigencharacter $\mathfrak h$. We denote the absolute value of this eigencharacter by $\eta$. By Remark \ref{rem:length}, for any $D\in \mbb N^{\Cal D}$ with $\varrho_X(D)=\ch\lambda$, we have $ \len (D) = \left< \mathfrak h + 2\rho_G, \ch\lambda\right>$. The effect of multiplying by $  (\eta\delta)^\frac{1}{2}(t)$ will be to replace $\ol\bbQ_\ell(-\dim \msf b)$, in the expression above, by $\ol\bbQ_\ell(\frac{1}{2})$, for those $\msf b$ in $\sY_{X^\bullet}^{D}$; this proves Theorem \ref{theorem-affine-closure}(i),(iv),(v).  

For the rest of the components $\msf b$, taking into account that $\ch\lambda\ge \ch\theta$ and $\mf h$ is trivial on the roots, therefore $\left<\mathfrak h ,\ch\lambda \right> = \left<\mathfrak h, \ch\theta\right>$ and $\left< 2\rho_G, \ch\lambda\right>$ has the same parity as $\left<  2\rho_G, \ch\theta\right>$, the effect of multiplying by $  (\eta\delta)^\frac{1}{2}(t)$ will be to replace $\ol\bbQ_\ell(-\dim \msf b)$ by $\ol\bbQ_\ell(\frac{\left<\mathfrak h + 2\rho_G,\ch\theta\right>}{2})$, proving Theorem \ref{theorem-general}. 

The remaining parts, (ii), (iii) of Theorem \ref{theorem-affine-closure} follow directly from Theorem \ref{thm:crystal-properties}.
\end{proof}




\subsection{Asymptotics and the basic function} \label{sect:asymptotics}
\source{intersection-complexes-unramified-l-factors}

We explain the proof of Corollary \ref{cor:asymptotics}, computing the basic function and its asymptotics.

\begin{proof}[Proof of Corollary \ref{cor:asymptotics}]
The (function-theoretic) Radon transform $\pi_{\emptyset !}$ on $X_\emptyset^\bullet$ is well-understood; in particular, the function $1_{X_\emptyset^\bullet(\mathfrak o)}$ maps to 
\[ \delta^{-\frac{1}{2}} \tr_{\check T}(\Fr, \Sym^\bullet (\check{\mathfrak n})) \cdot \tr_{\check T}(\Fr, \Sym^\bullet (\check{\mathfrak n}(1)))^{-1},\]
in the notation of \eqref{equation-pushforward-IC}. Moreover, $\pi_{\emptyset !}$ is equivariant for the action of the torus $T(F)$. This implies the formula \eqref{equation-asymptotics-IC} for $e_\emptyset^*\Phi_0$.

This, in turn, implies the formula \eqref{equation-basicfunction} for $\Phi_0$, by \cite[Corollary 5.5]{SaSatake}, where it is proven that, for an appropriate parametrization of the $G(\mf o)$-orbits, any $G(\mf o)$-invariant function $\Phi$ on $X^\bullet(F)$ is equal to its asymptotic $e_\emptyset^* \Phi$, restricted to the antidominant lattice $\ch\Lambda^{-,\Fr}$. Note that that paper was written under the assumption of $G$ being split (and $X$ being ``wavefront'', which is automatic in our setting), but the proof of the result that we are quoting is valid without this assumption. (For example, ``quasi-split'' is enough for the Casselman--Shalika method, which is the basis of that argument, and for the asymptotics theory developed in \cite{SV}.)
\end{proof}



\subsection{Euler factors of global integrals} \label{sect:nfold-explanation}
\source{intersection-complexes-unramified-l-factors}


Finally, to demonstrate how our results apply directly to compute Euler factors of global integrals of automorphic forms, let us return to the setting of Examples \ref{example-nfold}, \ref{example-nfold-global}, in order to discuss the Euler factorization of the pertinent global period integral. This will not directly invoke the nearby cycles functor (hence, is based on results of previous sections only), but it makes use of the asymptotics map $e_\emptyset^*$ that we just recalled.

To recall, the group is $G=(\mathbb G_m\times \SL_2^n)/\mu_2$, and $X$ is the affine closure of $H\backslash G$, where $H$ is the product of the group $H_0$ of \eqref{example-nfold} with a copy of $\mathbb G_m$ embedded diagonally into $G$ as 
\begin{equation}\label{Gmembedding} a \mod \mu_2 \mapsto \left( a, \begin{pmatrix} a \\ & a^{-1}\end{pmatrix}^n \right ) \mod \mu_2.
\source{intersection-complexes-unramified-l-factors}
\end{equation}
We let $\Phi = \prod_v \Phi_v$ be a function on $X^\bullet(\mathbb A)$ as in Example \ref{example-nfold-global}.

Let $\phi\in \pi$ be a cusp form, where the restriction of $\pi$ to $\mathbb G_m$ is a character of the form $\chi_0 |\bullet|^s$ for $\chi_0$ unitary and some $s \gg 0$. Fix a nontrivial character $\psi$ of $\mbb A/\Bbbk$, identified as a character of the quotient $N^-/H_0$, where $N^-$ is the lower triangular unipotent subgroup (mapping to the additive group through summation of the entries). The standard method of writing $\phi$ in terms of its Whittaker function $W_\phi$ with respect to $(N^-,\psi)$ shows that the integral   
\[ \int_{G(\Bbbk)\backslash G(\mbb A)} \phi(g) \sum_{\gamma\in X^\bullet(\Bbbk)} \Phi(\gamma g) dg = \int_{X^\bullet(\mbb A)} \Phi(g) \int_{H(\Bbbk)\backslash H(\mbb A)} \phi(hg) dh dg\]
``unfolds'' to the Eulerian integral
 \[ \int_{H_0\backslash G (\mbb A)} W_\phi(g) \Phi(g) dg.\]
 
We can write the Whittaker function as $W_\Phi(g) = \prod_v W_{\phi, v}(g_v)$, so that $W_{\phi,v}(1) = 1$ for almost all $v$, and we can factorize the invariant measure on $H_0\backslash G$ so that for almost all $v$, $\mu_v(H_0\backslash G(\mf o_v)) = 1$. 

\begin{prop}
\source{intersection-complexes-unramified-l-factors}
\notready\label{prop:localzeta}
\source{intersection-complexes-unramified-l-factors}
At places $v$ where $W_{\phi,v}(1)$ is $G(\mf o_v)$-invariant with $W_{\phi,v}(1) = 1$ (in particular, $\pi_v$ is unramified), $\Phi_v=\Phi_{0,v}$ (the IC function), $\mu_v(H_0\backslash G(\mf o_v)) = 1$, and the conductor of $\psi|_{\Bbbk_v}$ is the ring of integers, we have
\begin{equation}\label{eq:localzetafactor}
\source{intersection-complexes-unramified-l-factors}
   \int_{H_0\backslash G (\Bbbk_v)} W_{\phi,v}(g_v) \Phi_v(g_v) dg_v = L(\pi_v, \otimes, 1 - \frac{n}{2}).
\end{equation}
\end{prop}

\begin{proof}
The proof can be obtained by direct computation from the explicit formula \eqref{equation-basicfunction} for the basic function. However, we would like to sketch a ``pure thought'' argument, which applies to other cases as well, without providing all details. 
We drop the index $v$, denoting the place under consideration, and write $F$ for $\Bbbk_v$.

First of all, the \emph{square of the absolute value} of \eqref{eq:localzetafactor} follows immediately from Plancherel-theoretic considerations. Indeed, we can write the local integral as 
 \[    Z(\Phi, \pi)=  \int_{N^-\backslash G (F)} W_{\phi}(g) W_{\Phi}(g)  dg,\]
 where 
 \begin{equation}\label{unfolding} W_{\Phi}(g) = \int_{H_0\backslash N^-(F)} \Phi(ng)\psi(n) dn.
\source{intersection-complexes-unramified-l-factors}
 \end{equation}
 We identify the abelianization $G^{\text{ab}}$ with $\mathbb G_m$, so that the composite $\mathbb G_m\to G\to\mathbb G_m$ is the square character. 
 Let $\eta$ be the character $a\mapsto |a|^{1-n}$ on $G$ --- it is the character by which it acts on the $\SL_2^n$-invariant measure on $X^\bullet$. Fix that measure giving volume $1$ to $X^\bullet(\mf o)$.
 It is known from \cite[Theorem 9.5.9]{SV} that the ``unfolding'' map $\Phi \mapsto W_{\Phi}$ extends to an $L^2$-isometry
 \[ L^2(X^\bullet(F)) \xrightarrow\sim L^2(N^-\backslash G(F), \psi^{-1}),\]
 where the Haar measure on $N^-\backslash G(F)$ is fixed so that the volume of $N^-\backslash G(\mf o)$ is $1$.
 
 The local zeta integrals appear in the Plancherel decomposition of Whittaker functions,  \begin{equation}\label{eq:PlancherelWhittaker} \Vert \Phi\Vert^2_{L^2(X^\bullet(F))} = \Vert W_{\Phi}\Vert^2_{L^2(N(F)\backslash G(F), \psi^{-1})} = \int |Z(\Phi, \pi)|^2 d\pi, 
\source{intersection-complexes-unramified-l-factors}
 \end{equation}
  where $\pi$ ranges over the unitary unramified dual of $G(F)$, which can be identified with the set of semisimple conjugacy classes in the compact form of the complex dual group $\check G(\mbb C)$, and the Plancherel measure $d\pi$ is the Weyl measure $\left|(1-e^{\ch\alpha})\right|^2 d\chi$ (when we identify the unramified dual with the quotient $\ch T^1/W$, where $\ch T^1$ is the group of unitary unramified characters --- the maximal compact subgroup of the complex points of the dual Cartan $\ch T$).
 
 On the other hand, the Plancherel decomposition of $\Phi$ can also be expressed in terms of its asymptotics:
 \[ \Vert \Phi\Vert^2_{L^2(X^\bullet(F))} = \frac{1}{|W|} \Vert e_\emptyset^* \Phi\Vert^2_{L^2(X_\emptyset^\bullet(F))},
 \]
 where we have applied \cite[Theorem 7.3.1]{SV}, together with the following observations: 1) our calculation of $e_\emptyset^* \Phi$ shows that it already lies in $L^2(X_\emptyset^\bullet(F))$, so there is no difference between that and what is denoted by $\iota_\emptyset^* \Phi$ in \emph{loc.cit.}; 2) there are no unramified functions $\Phi$ on $X^\bullet(F)$ with $e_\emptyset^*\Phi = 0$; this is essentially \cite[Theorem 6.2.1]{SaSpc}, and it implies, together with the previous point, that the summands of \cite[Theorem 7.3.1]{SV} with $\Theta\ne \emptyset$ do not contribute to the Plancherel decomposition. (This is not an essential point for this proof; the argument would go through even if there were contributions outside of the most continuous spectrum.)

 
By our formula \eqref{equation-pushforward-IC-expl}, as specialized in Example \ref{example-nfold}, and the commutation of asymptotics with Radon transform, \eqref{Radon-asymptotics}, we can write the Plancherel decomposition of the basic function as 
 \begin{equation}\label{eq:Plancherelnfold} \Vert \Phi\Vert^2_{L^2(X^\bullet(F))} = \frac{1}{|W|} \int_{\check T^1} 
\source{intersection-complexes-unramified-l-factors}
\left|\frac{\prod_{\check\alpha\in\check\Phi^+} (1-e^{\check\alpha}(\chi))}{\prod_{\check\lambda \in \mB^+} (1-q^{-\frac{1}{2}} e^{\check\lambda}(\chi))}\right|^2 d\chi =   
 \int L(\pi, \otimes, \frac{1}{2}) L(\tilde \pi, \otimes, \frac{1}{2})
 d\pi.
 \end{equation}
 
Comparing \eqref{eq:PlancherelWhittaker} and \eqref{eq:Plancherelnfold}, we get that $|Z(\Phi, \pi)|^2 = L(\pi, \otimes, \frac{1}{2}) L(\tilde\pi, \otimes, \frac{1}{2})$ for $\pi$ unitary. This is what we get ``for free'' from the $L^2$-decomposition.
 
To upgrade it to a formula on $Z(\Phi, \pi)$, we utilize the following information about the image $W_{\Phi}$ under the unfolding map. Note that this is a $G(\mf o)$-invariant Whittaker function, hence can be identified as a function on the \emph{antidominant} coweights, by evaluating it on the elements $t^{\ch\theta}$. The facts that we need are:

\begin{enumerate}
 \item The support of $W_{\Phi}$ lies in the intersection of the antidominant lattice with the cone generated by $\mf c_X^-$ and the coweight \[\ch\theta_0:= - \frac{\check\alpha_1+\check\alpha_2+\dots+\check\alpha_n - \check m}{2}.\]
 Its value at $1 = t^0$ is $1$.
 
 Indeed, the support of $\Phi$ lies in $H \xi t^{\mf c_X^-} G(\mf o)$, where $\xi$ is the diagonal image of $\begin{pmatrix} 1 & 0 \\ 1 & 1\end{pmatrix}$, and from the definition \eqref{unfolding} of the unfolding map, the support of $W_{\Phi}$ will lie in $H \xi N^- t^{\mf c_X^-}  G(\mf o) = N^- \mbb G_m t^{\mf c_X^-}  G(\mf o)$, where $\mbb G_m$ is embedded as in \eqref{Gmembedding}. But a $G(\mf o)$-invariant Whittaker function with respect to a character of $N^-$ whose conductor is $\mf o$ can only be supported on antidominant elements. 
 
 When $g=1$, the integrand of \eqref{unfolding} only lies in the support of $\Phi$ when $n\in H_0\backslash N^-(\mf o)$.
 
 \item The value of $W_{\Phi}$ at any $t^{\ch\theta}$ is a polynomial in $q^{-\frac{1}{2}}$. 
 
 This fact, which can be seen as expressing the ``motivic'' nature of this function, requires some explanation. First of all, the statement is true if $W_{\Phi}$ is replaced by the pushforward $\pi_! \Phi$, by Theorem \ref{theorem-affine-closure} (see also \S \ref{sect:finitefields}). Secondly, the unfolding integral \eqref{unfolding} of the characteristic function of each $G(\mf o)$-orbit has this property; this can be seen by direct calculation, or by a similar geometric argument, and we omit the details. 
 
 \item If $\varphi_{\ch\theta}$ denotes the $G(\mf o)$-invariant Whittaker function supported on the coset $N^-t^{\ch\theta} G(\mf o)$ (with $\ch\theta$ antidominant) and equal to $1$ on $\ch\theta$, its asymptotic $e_\emptyset^* \varphi_{\ch\theta}$ is supported on the union of a finite number of cosets $N^-t^{\ch\theta'} G(\mf o)$, with $\ch\theta'-\ch\theta$ in the positive root monoid, and is a polynomial in $q^{-\frac{1}{2}}$. 
 
 Here, the asymptotics map is for the Whittaker model, but we denote it by the same symbol. It takes values in $N^-\backslash G(F)$, without a character on $N^-$. This fact is a corollary of the Casselman--Shalika formula. (See \cite[Theorem 6.8 and Example 6.4]{SaSatake} for an intepretation of the Casseman--Shalika formula in terms of asymptotics.)
 

\end{enumerate}

Granted, now, the facts above, we can represent $e_\emptyset^* W_{\Phi}$ as a formal series in the elements $e^{\ch\theta}$, which stand for the characteristic functions of the sets $N^- t^{\ch\theta} G(\mf o)$, multiplied by $\delta^{-\frac{1}{2}}(t^{\ch\theta}) = q^{\left<\rho_G,\ch\theta\right>}$, with coefficients in $\mathbb C[q^{-\frac{1}{2}}]$:
\[ e_\emptyset^* W_{\Phi} = \sum_{i,\ch\theta} c_{i,\ch\theta} q^{-\frac{i}{2}} e^{\ch\theta}.\]

According to the first and the third facts above, the support of this function lies in the set of $\ch\theta$'s which are translates, by the positive coroot monoid, of the intersection
\[ (\mf c_X^- + \text{span}(\ch\theta_0)) \cap \ch\Lambda^-.\]
The description of colors of $X$ in Example \ref{example-nfold} allows us to conclude that this monoid only includes coweights for which the ``determinant'' character 
\[\frac{\kappa \ch\mu + \sum_i \kappa_i \ch\alpha_i}{2} \mapsto \kappa\]
is \emph{nonnegative}; moreover, the restriction to the kernel of the determinant
cocharacter is simply equal to the asymptotics of the ``basic Whittaker function'' $\varphi_0$, which by \cite[Theorem 6.8 and Example 6.4]{SaSatake} is equal to 
\[\prod_{\ch\alpha \in \ch\Phi^+} (1-e^{\ch\alpha}).\]

Thus, if we ``divide'' the function $e_\emptyset^* W_{\Phi}$ by the product above (this corresponds to acting on it by a series in the Hecke algebra of the torus $T$), we obtain another series
\[\prod_{\ch\alpha \in \ch\Phi^+} (1-e^{\ch\alpha})^{-1} e_\emptyset^* W_{\Phi} = 
1\cdot e^0 + \sum_{j=1}^\infty \sum_{i=0}^\infty q^{-\frac{i}{2}} e^{j\frac{\ch\mu}{2}} C_{i,j}
,\]
where the $C_{i,j}$'s are finite linear combinations of the elements $e^{\ch\lambda}$, where $\ch\lambda$ now ranges over half-multiples of elements in the coroot lattice. (The finiteness of the linear combination also follows from \cite[Theorem 6.8 and Example 6.4]{SaSatake}: the asymptotics of \emph{every} $\varphi_{\ch\theta}$ is a ``multiple'' of the factor $\prod_{\ch\alpha \in \ch\Phi^+} (1-e^{\ch\alpha})$.)
Note that $\frac{\ch\mu}{2}$ is not, by itself, a coweight of $G$, so one has to expand this series to make sense of it as a linear combination of the elements $e^{\ch\theta}$ with $\ch\theta \in \ch\Lambda$.

Now we invoke the Plancherel formula: the fact that the unfolding map is an isometry tells us that the Plancherel density of $W_{\Phi}$ is also given by \eqref{eq:Plancherelnfold}. On the other hand, this Plancherel density can be expressed in terms of the asymptotics
$e_\emptyset^* W_{\Phi}$ as 
\[ \Vert W_{\Phi}\Vert^2 = \frac{1}{|W|} \Vert e_\emptyset^* W_{\Phi}\Vert^2 = \frac{1}{|W|} \int_{\check T^1} 
\left|1\cdot e^0 + \sum_{j=1}^\infty \sum_{i=0}^\infty q^{-\frac{i}{2}} e^{j\frac{\ch\mu}{2}} C_{i,j}\right|^2 \left| \prod_{\ch\alpha \in \ch\Phi^+} (1-e^{\ch\alpha})\right|^2 d\chi = \]
\[ = \frac{1}{|W|} \int_{\check T^1} 
\left(1\cdot e^0 + \sum_{j=1}^\infty \sum_{i=0}^\infty q^{-\frac{i}{2}} e^{j\frac{\ch\mu}{2}} C_{i,j}\right) \left(1\cdot e^0 + \sum_{j=1}^\infty \sum_{i=0}^\infty q^{-\frac{i}{2}} e^{-j\frac{\ch\mu}{2}} C^*_{i,j}\right) \prod_{\ch\alpha \in \ch\Phi^+} \left|(1-e^{\ch\alpha})\right|^2 d\chi,\]
where we now identify the elements $e^{\ch\theta}$ with characters of the dual torus $\ch T$, use the fact that for a unitary character $\overline{e^{\ch\theta}(\chi)} = e^{-\ch\theta}(\chi)$, and set $C^*_{i,j} = \sum_{\ch\lambda} \overline{c_{\ch\lambda}} e^{-\ch\lambda}$ if $C_{i,j} = \sum_{\ch\lambda} c_{\ch\lambda} e^{\ch\lambda} $. 

In other words, we have expressed the Plancherel density 
\[ 
 L(\pi, \otimes, \frac{1}{2}) L(\tilde \pi, \otimes, \frac{1}{2})
\]
as the product 
\[
 \left(1\cdot e^0 + \sum_{j=1}^\infty \sum_{i=0}^\infty q^{-\frac{i}{2}} e^{j\frac{\ch\mu}{2}} C_{i,j}\right) \left(1\cdot e^0 + \sum_{j=1}^\infty \sum_{i=0}^\infty q^{-\frac{i}{2}} e^{-j\frac{\ch\mu}{2}} C^*_{i,j}\right). 
\]
We will be done if we can identify the first factor of the former with the first factor of the latter. But, viewed as series in the elements $q^{-\frac{i}{2}} e^{-j\frac{\ch\mu}{2}}$ (with coefficients in polynomials in the group ring of the half-root lattice), the first factor of the former is the restriction of the series to the elements of the form $q^{-\frac{i}{2}} e^{i\frac{\ch\mu}{2}}$, while the product 
$ L(\pi, \otimes, \frac{1}{2}) L(\tilde \pi, \otimes, \frac{1}{2})$ is supported on elements of the form $q^{-\frac{i}{2}} e^{j\frac{\ch\mu}{2}}$ with $j\le i$. It follows that the series $ 1\cdot e^0 + \sum_{j=1}^\infty \sum_{i=0}^\infty q^{-\frac{i}{2}} e^{j\frac{\ch\mu}{2}} C_{i,j}$ is also supported on elements of the form $q^{-\frac{i}{2}} e^{j\frac{\ch\mu}{2}}$ with $j \le i$, coincides with $L(\pi, \otimes, \frac{1}{2})$ on elements with $i=j$, and a simple inductive argument in the variable $i-j$ shows that it coincides with it everywhere.
\end{proof}
 


\appendix
